\section{The boundary value of the definite difference}

% BEGIN B75S1
\paragraph{B75S1. Removing the first column from the theta series.}
Let $E$ be the CM field, $\Psi$ its CM type, $r\ge1$ the number of
columns and $n$ the number of rows. Write $e$ for the first standard
column vector. Given a period tuple $Z$ with $Y_\tau=\operatorname{Im}
Z_\tau>0$ and $v\in E^r$, put
\[
   Z_t(v)_\tau=Z_\tau+it\,v_\tau v_\tau^*,\qquad
   v_\tau=\tau(v),\quad t\ge0.
\]
Every such tuple is in the half space. Write $Z'$ for the principal
submatrix obtained by deleting the first row and column at every place;
its imaginary part is positive definite. For an original Hermitian
Schwartz--Bruhat test $f$ put $f_0(u)=f([0\mid u])$. Its denominator
bound and periods restrict to the remaining columns, so $f_0$ is again
Schwartz--Bruhat. For any fixed pair of elliptic variables $z,w$ let
$z',w'$ denote those same variables without their first columns. Then
\[
   \lim_{t\to\infty}\Theta_r(f;Z_t(e),z,w)
           =\Theta_{r-1}(f_0;Z',z',w').
\]

Indeed the summand at $x$ is multiplied by
$\exp(-2\pi t\sum_{\tau,k}|\tau(x_{k1})|^2)$. The support lies in a
fixed denominator lattice in Minkowski space, and $f$ is bounded there
by periodicity. At the fixed $Z,z,w$ the terms in absolute value are
bounded by a constant times
$\exp(-a\|x\|^2+b\|x\|)$ on this lattice for $a>0$ and some $b\ge0$.
This is summable and still dominates after the extra factor is inserted.
A CM type contains an embedding, so the limit of the factor is zero
unless the first column vanishes. For $x=[0\mid u]$ the original
quadratic and both linear exponents are exactly the smaller ones.
Dominated convergence proves the formula, including $r=1$.
% END B75S1

% BEGIN B75S2
\paragraph{B75S2. Surviving bottom spaces and their finite tests.}
For the Eisenstein series suppose $n$ is even and $n>4r$. Write a
rational split unitary matrix $g$ in blocks $(A,B;C,D)$ and write
$U$ for its bottom row space. The condition $Ce=0$ means every row
in $U$ is perpendicular for the split form to the vector
$\ell=(0,e^t)$ in the second half. Since $U$ is maximal isotropic it
equals its perpendicular. Thus $Ce=0$ is intrinsic and is equivalent
to $E\ell\subset U$. Passage to $\ell^\perp/E\ell$, identified
with the split space of degree $r-1$ by deleting both distinguished
coordinates, bijects these $U$ with the maximal isotropic spaces $U'$
of the smaller degree. The inverse takes the entire preimage in
$\ell^\perp$. Every $U'$ has a rational split unitary representative
$g'$ by hyperbolic completion as in C75E1 (the empty matrix suffices
in degree zero). Extending $g'$ by the identity on the distinguished
hyperbolic plane gives a representative $\iota g'$ for this preimage.

For every $Z$, including periods with mixed entries in its first row
and column, one has at each embedding
\[
     \det(C_{\iota g'}Z+D_{\iota g'})=\det(C_{g'}Z'+D_{g'}).
\]
Indeed the first row of the larger denominator is $(1,0,\ldots,0)$
and its complementary minor is the displayed lower denominator.
The finite tests satisfy
\[
      (\omega_r(\iota g')f)_0=\omega_{r-1}(g')f_0 .
\]
To see this, at each place insert the period
$\operatorname{diag}(i(1+t),Z')$ and elliptic arguments with first
columns zero into the defining theta transformation. Block
multiplication shows that its transformed first period is still
$i(1+t)$, the remaining period and elliptic variables transform by
$g'$, and the two determinant and Gaussian factors together give
exactly the smaller factor of automorphy. Apply B75S1 on both sides
as $t\to\infty$. The resulting equality for every pair of smaller
elliptic variables is the defining theta transformation for the
restricted test. For clarity about uniqueness, at a fixed smaller
period restrict the two elliptic variables to conjugate real-slice
pairs. The linear exponents are then the real trace pairings with
the embedded rational matrices. On a common period torus for their
denominator lattice these are distinct characters. The Gaussian at
every frequency is nonzero and B75S1 gives absolute summability,
so integrating against one character recovers its test value. The
linear elliptic changes at $g'$ are invertible before taking this
slice. In degree zero there is a single coefficient. This proves the
finite identity and, by evaluating at zero, the required equality of
the two Eisenstein numerators.
% END B75S2

% BEGIN B75S3
\paragraph{B75S3. The Eisenstein limit at zero elliptic arguments.}
With the data of B75S2,
\[
    \lim_{t\to\infty}\operatorname{Eis}_r(f;Z_t(e),0,0)
          =\operatorname{Eis}_{r-1}(f_0;Z',0,0).
\]
At one embedding put $Q=CZ+D$ and $H=Q^{-1}C$. The denominator
is invertible, and $-\operatorname{Im}H=HYH^*$ by O74E4. For
any column $u$ the rank-one determinant formula gives
\[
  |\det(Q+itCuu^*)|
    =|\det Q|\,|1+it u^*Hu|
    \ge |\det Q|\,(1+t\,u^*HYH^*u).
\]
The last quadratic form is positive if $Cu\ne0$. Indeed set
$b=(Q^*)^{-1}u$. If $H^*u=C^*b=0$, then
$u=Q^*b=D^*b$ and $Cu=CD^*b=DC^*b=0$ by isotropy.
Consequently each inverse denominator power at $Z_t(e)$ has absolute
value at most the one at $Z$. If $Ce\ne0$, it tends to zero at each
embedding, since $C,e$ are rational over $E$ and $n>0$.
The Weil numerator is constant in $t$. The absolute sum at $Z$ in
C75E4 is therefore a summable majorant for the entire ray.

For the survivors, the parabolic calculation of O74E3 applies to any
test: changing representatives with the same bottom space multiplies
both numerator and automorphy factor by the same scalar. Use the
representatives $\iota g'$ of B75S2. Their denominators are constant
along this ray, and B75S2 identifies each whole term with the smaller
one. The bijection of bottom spaces has multiplicity one, so dominated
convergence gives the stated limit.
% END B75S3

% BEGIN B75S4
\paragraph{B75S4. Every rational direction with the original admission.}
Fix $M,R>0$ and an admitted family $\phi=(\phi_\Lambda)_\Lambda$
in degree $r$. Let $v\ne0$ belong to $E^r$. Choose $A\in
\mathrm{GL}_r(E)$ with $Av=e$ by extending $v$ to a basis. Put
\[
 g=\operatorname{diag}(A,(A^{-1})^*),\qquad
 j=\prod_{\tau\in\Psi}\tau(\det((A^{-1})^*))^n\ne0,\qquad
 \psi_\Lambda(x)=j\phi_\Lambda(xA).
\]
The native Levi transformation gives $\psi_\Lambda=\omega(g)\phi_\Lambda$:
in the theta formula the new summation index is $xA$, and the
upper-right block is zero, so the parabolic quadratic character is one.

Choose $k>0$ clearing the denominators of both $A$ and $A^{-1}$ and
put $R'=kR$. We check admission at $(M,R')$. Positivity is immediate
and the definite row action commutes with right multiplication by
$A$. Whenever $\psi_\Lambda(x)\ne0$, $RxA$ has columns in $\Lambda$,
so $R'x=(RxA)(kA^{-1})$ does also. Adding $R'y$ with columns of $y$
in $\Lambda$ changes $xA$ by $R(y\,kA)$, an original period.
For locality suppose $\Lambda,\Lambda'$ agree at $MR'$ and $R'x$
has columns in both. An integer $s>0$ prime to $MR'$ carries each
lattice into the other by multiplication. Set $u=RxA$. The columns
of $k^2u=(R'x)(kA)$ belong to both lattices. If $u$ has columns in
one, $su$ has columns in the other. Since $\gcd(s,k^2)=1$, B\'ezout
then places $u$ in the other as well. If neither contains the columns
of $u$, both original values vanish by support. Otherwise locality
at $MR$ applies, using the same $s$. This proves equality of the new
values. Finally $\lambda_\Lambda(y)=\psi_\Lambda([0\mid y])$ retains
all five clauses at $(M,R')$. In the support and locality tests a
first column of zero belongs to both lattices, and padding a period
with a zero first column gives a period of the larger test.

Write $W_\tau=A_\tau Z_\tau A_\tau^*$ and $W'$ for its first
principal deletion. Then $gZ_t(v)=W_t(e)$. At zero elliptic
arguments the theta transformation multiplies by $j$. The same holds
for Eisenstein: right multiplication by $g$ bijects the rational
bottom spaces, the Weil composition law changes the new numerator
to $(\omega(\sigma g)\phi_\Lambda)(0)$, and the denominator cocycle
gives $j(\sigma,gZ)=j(\sigma g,Z)/j$. Representative independence
and C75E4 justify reindexing. Thus the individual theta and Eisenstein
limits in the $v$ direction are respectively $j^{-1}$ times the
smaller theta and Eisenstein values of $\lambda_\Lambda$ at $W',0,0$.

Let $D_{r,M,\phi}$ be the native weighted sum over the finite definite
class set of these theta values minus the Eisenstein values at zero
elliptic arguments. The class set, its chosen lattices and its weights
depend on the row problem and $M$; they are identical in both column
degrees. Taking the finite sum gives
\[
   \lim_{t\to\infty}D_{r,M,\phi}(Z_t(v))
           =j^{-1}D_{r-1,M,\lambda}(W').
\]
The assertion $\mathrm{SWAt}(n,r-1)$ quantifies over every positive
level and every admitted lower-degree test. It therefore makes the
right side zero, including after the change to $R'$. This proves
the rational cusp limit needed in the degree induction.
% END B75S4
